%% pipeline_overview.tex - TikZ diagram of LLM Evaluation Pipeline

% This file defines a figure that can be included in main.tex
% Use: %% pipeline_overview.tex - TikZ diagram of LLM Evaluation Pipeline

% This file defines a figure that can be included in main.tex
% Use: %% pipeline_overview.tex - TikZ diagram of LLM Evaluation Pipeline

% This file defines a figure that can be included in main.tex
% Use: %% pipeline_overview.tex - TikZ diagram of LLM Evaluation Pipeline

% This file defines a figure that can be included in main.tex
% Use: \input{figures/pipeline_overview}

\begin{figure}[htbp]
\centering
\resizebox{0.95\textwidth}{!}{%
\begin{tikzpicture}[
    node distance=1.0cm and 1.2cm,
    box/.style={
        rectangle,
        draw=blue!70!black,
        fill=blue!10,
        rounded corners,
        minimum height=0.9cm,
        minimum width=2cm,
        text centered,
        font=\small,
        align=center
    },
    databox/.style={
        rectangle,
        draw=green!60!black,
        fill=green!10,
        rounded corners,
        minimum height=0.7cm,
        minimum width=1.8cm,
        text centered,
        font=\small,
        align=center
    },
    evalbox/.style={
        rectangle,
        draw=orange!70!black,
        fill=orange!10,
        rounded corners,
        minimum height=0.9cm,
        minimum width=2cm,
        text centered,
        font=\small,
        align=center
    },
    outputbox/.style={
        rectangle,
        draw=purple!70!black,
        fill=purple!10,
        rounded corners,
        minimum height=0.9cm,
        minimum width=2cm,
        text centered,
        font=\small,
        align=center
    },
    arrow/.style={
        ->,
        >=stealth,
        thick,
        draw=gray!70
    }
]

% Test Set
\node[databox] (testset) {Test Set};

% Optional: Retrieved Context for RAG
\node[databox, below=0.6cm of testset] (context) {Retrieved Context};

% LLM Application
\node[box, right=1.2cm of testset] (llm) {LLM Application};

% Output
\node[box, right=1.2cm of llm] (output) {Generated Output};

% Scoring Methods (branching)
\node[evalbox, above right=0.4cm and 1.2cm of output] (automated) {Automated Metrics};
\node[evalbox, right=1.2cm of output] (llmjudge) {LLM-as-Judge};
\node[evalbox, below right=0.4cm and 1.2cm of output] (human) {Human Eval};

% Aggregation
\node[outputbox, right=1.2cm of llmjudge] (scores) {Quality Scores};

% Dashboard/Reporting
\node[outputbox, right=1.0cm of scores] (dashboard) {Dashboard};

% Arrows
\draw[arrow] (testset) -- (llm);
\draw[arrow] (context) -| (llm);
\draw[arrow] (llm) -- (output);
\draw[arrow] (output) -- (automated);
\draw[arrow] (output) -- (llmjudge);
\draw[arrow] (output) -- (human);
\draw[arrow] (automated) -- (scores);
\draw[arrow] (llmjudge) -- (scores);
\draw[arrow] (human) -- (scores);
\draw[arrow] (scores) -- (dashboard);

% Feedback loop
\draw[arrow, dashed, bend left=40] (dashboard.south) to node[below, font=\footnotesize, text=gray] {Iterate} (testset.south);

% Labels
\node[above=0.1cm of testset, font=\footnotesize\bfseries, text=green!60!black] {Inputs};
\node[above=0.1cm of llm, font=\footnotesize\bfseries, text=blue!70!black] {System};
\node[above=0.1cm of llmjudge, font=\footnotesize\bfseries, text=orange!70!black] {Evaluation};
\node[above=0.1cm of dashboard, font=\footnotesize\bfseries, text=purple!70!black] {Outputs};

\end{tikzpicture}%
}
\caption{LLM Evaluation Pipeline Overview. Test inputs flow through the LLM application to generate outputs, which are scored by automated metrics, LLM judges, and/or human evaluators. Aggregated scores feed dashboards and inform iterative improvements.}
\label{fig:pipeline-overview}
\end{figure}


\begin{figure}[htbp]
\centering
\resizebox{0.95\textwidth}{!}{%
\begin{tikzpicture}[
    node distance=1.0cm and 1.2cm,
    box/.style={
        rectangle,
        draw=blue!70!black,
        fill=blue!10,
        rounded corners,
        minimum height=0.9cm,
        minimum width=2cm,
        text centered,
        font=\small,
        align=center
    },
    databox/.style={
        rectangle,
        draw=green!60!black,
        fill=green!10,
        rounded corners,
        minimum height=0.7cm,
        minimum width=1.8cm,
        text centered,
        font=\small,
        align=center
    },
    evalbox/.style={
        rectangle,
        draw=orange!70!black,
        fill=orange!10,
        rounded corners,
        minimum height=0.9cm,
        minimum width=2cm,
        text centered,
        font=\small,
        align=center
    },
    outputbox/.style={
        rectangle,
        draw=purple!70!black,
        fill=purple!10,
        rounded corners,
        minimum height=0.9cm,
        minimum width=2cm,
        text centered,
        font=\small,
        align=center
    },
    arrow/.style={
        ->,
        >=stealth,
        thick,
        draw=gray!70
    }
]

% Test Set
\node[databox] (testset) {Test Set};

% Optional: Retrieved Context for RAG
\node[databox, below=0.6cm of testset] (context) {Retrieved Context};

% LLM Application
\node[box, right=1.2cm of testset] (llm) {LLM Application};

% Output
\node[box, right=1.2cm of llm] (output) {Generated Output};

% Scoring Methods (branching)
\node[evalbox, above right=0.4cm and 1.2cm of output] (automated) {Automated Metrics};
\node[evalbox, right=1.2cm of output] (llmjudge) {LLM-as-Judge};
\node[evalbox, below right=0.4cm and 1.2cm of output] (human) {Human Eval};

% Aggregation
\node[outputbox, right=1.2cm of llmjudge] (scores) {Quality Scores};

% Dashboard/Reporting
\node[outputbox, right=1.0cm of scores] (dashboard) {Dashboard};

% Arrows
\draw[arrow] (testset) -- (llm);
\draw[arrow] (context) -| (llm);
\draw[arrow] (llm) -- (output);
\draw[arrow] (output) -- (automated);
\draw[arrow] (output) -- (llmjudge);
\draw[arrow] (output) -- (human);
\draw[arrow] (automated) -- (scores);
\draw[arrow] (llmjudge) -- (scores);
\draw[arrow] (human) -- (scores);
\draw[arrow] (scores) -- (dashboard);

% Feedback loop
\draw[arrow, dashed, bend left=40] (dashboard.south) to node[below, font=\footnotesize, text=gray] {Iterate} (testset.south);

% Labels
\node[above=0.1cm of testset, font=\footnotesize\bfseries, text=green!60!black] {Inputs};
\node[above=0.1cm of llm, font=\footnotesize\bfseries, text=blue!70!black] {System};
\node[above=0.1cm of llmjudge, font=\footnotesize\bfseries, text=orange!70!black] {Evaluation};
\node[above=0.1cm of dashboard, font=\footnotesize\bfseries, text=purple!70!black] {Outputs};

\end{tikzpicture}%
}
\caption{LLM Evaluation Pipeline Overview. Test inputs flow through the LLM application to generate outputs, which are scored by automated metrics, LLM judges, and/or human evaluators. Aggregated scores feed dashboards and inform iterative improvements.}
\label{fig:pipeline-overview}
\end{figure}


\begin{figure}[htbp]
\centering
\resizebox{0.95\textwidth}{!}{%
\begin{tikzpicture}[
    node distance=1.0cm and 1.2cm,
    box/.style={
        rectangle,
        draw=blue!70!black,
        fill=blue!10,
        rounded corners,
        minimum height=0.9cm,
        minimum width=2cm,
        text centered,
        font=\small,
        align=center
    },
    databox/.style={
        rectangle,
        draw=green!60!black,
        fill=green!10,
        rounded corners,
        minimum height=0.7cm,
        minimum width=1.8cm,
        text centered,
        font=\small,
        align=center
    },
    evalbox/.style={
        rectangle,
        draw=orange!70!black,
        fill=orange!10,
        rounded corners,
        minimum height=0.9cm,
        minimum width=2cm,
        text centered,
        font=\small,
        align=center
    },
    outputbox/.style={
        rectangle,
        draw=purple!70!black,
        fill=purple!10,
        rounded corners,
        minimum height=0.9cm,
        minimum width=2cm,
        text centered,
        font=\small,
        align=center
    },
    arrow/.style={
        ->,
        >=stealth,
        thick,
        draw=gray!70
    }
]

% Test Set
\node[databox] (testset) {Test Set};

% Optional: Retrieved Context for RAG
\node[databox, below=0.6cm of testset] (context) {Retrieved Context};

% LLM Application
\node[box, right=1.2cm of testset] (llm) {LLM Application};

% Output
\node[box, right=1.2cm of llm] (output) {Generated Output};

% Scoring Methods (branching)
\node[evalbox, above right=0.4cm and 1.2cm of output] (automated) {Automated Metrics};
\node[evalbox, right=1.2cm of output] (llmjudge) {LLM-as-Judge};
\node[evalbox, below right=0.4cm and 1.2cm of output] (human) {Human Eval};

% Aggregation
\node[outputbox, right=1.2cm of llmjudge] (scores) {Quality Scores};

% Dashboard/Reporting
\node[outputbox, right=1.0cm of scores] (dashboard) {Dashboard};

% Arrows
\draw[arrow] (testset) -- (llm);
\draw[arrow] (context) -| (llm);
\draw[arrow] (llm) -- (output);
\draw[arrow] (output) -- (automated);
\draw[arrow] (output) -- (llmjudge);
\draw[arrow] (output) -- (human);
\draw[arrow] (automated) -- (scores);
\draw[arrow] (llmjudge) -- (scores);
\draw[arrow] (human) -- (scores);
\draw[arrow] (scores) -- (dashboard);

% Feedback loop
\draw[arrow, dashed, bend left=40] (dashboard.south) to node[below, font=\footnotesize, text=gray] {Iterate} (testset.south);

% Labels
\node[above=0.1cm of testset, font=\footnotesize\bfseries, text=green!60!black] {Inputs};
\node[above=0.1cm of llm, font=\footnotesize\bfseries, text=blue!70!black] {System};
\node[above=0.1cm of llmjudge, font=\footnotesize\bfseries, text=orange!70!black] {Evaluation};
\node[above=0.1cm of dashboard, font=\footnotesize\bfseries, text=purple!70!black] {Outputs};

\end{tikzpicture}%
}
\caption{LLM Evaluation Pipeline Overview. Test inputs flow through the LLM application to generate outputs, which are scored by automated metrics, LLM judges, and/or human evaluators. Aggregated scores feed dashboards and inform iterative improvements.}
\label{fig:pipeline-overview}
\end{figure}


\begin{figure}[htbp]
\centering
\resizebox{0.95\textwidth}{!}{%
\begin{tikzpicture}[
    node distance=1.0cm and 1.2cm,
    box/.style={
        rectangle,
        draw=blue!70!black,
        fill=blue!10,
        rounded corners,
        minimum height=0.9cm,
        minimum width=2cm,
        text centered,
        font=\small,
        align=center
    },
    databox/.style={
        rectangle,
        draw=green!60!black,
        fill=green!10,
        rounded corners,
        minimum height=0.7cm,
        minimum width=1.8cm,
        text centered,
        font=\small,
        align=center
    },
    evalbox/.style={
        rectangle,
        draw=orange!70!black,
        fill=orange!10,
        rounded corners,
        minimum height=0.9cm,
        minimum width=2cm,
        text centered,
        font=\small,
        align=center
    },
    outputbox/.style={
        rectangle,
        draw=purple!70!black,
        fill=purple!10,
        rounded corners,
        minimum height=0.9cm,
        minimum width=2cm,
        text centered,
        font=\small,
        align=center
    },
    arrow/.style={
        ->,
        >=stealth,
        thick,
        draw=gray!70
    }
]

% Test Set
\node[databox] (testset) {Test Set};

% Optional: Retrieved Context for RAG
\node[databox, below=0.6cm of testset] (context) {Retrieved Context};

% LLM Application
\node[box, right=1.2cm of testset] (llm) {LLM Application};

% Output
\node[box, right=1.2cm of llm] (output) {Generated Output};

% Scoring Methods (branching)
\node[evalbox, above right=0.4cm and 1.2cm of output] (automated) {Automated Metrics};
\node[evalbox, right=1.2cm of output] (llmjudge) {LLM-as-Judge};
\node[evalbox, below right=0.4cm and 1.2cm of output] (human) {Human Eval};

% Aggregation
\node[outputbox, right=1.2cm of llmjudge] (scores) {Quality Scores};

% Dashboard/Reporting
\node[outputbox, right=1.0cm of scores] (dashboard) {Dashboard};

% Arrows
\draw[arrow] (testset) -- (llm);
\draw[arrow] (context) -| (llm);
\draw[arrow] (llm) -- (output);
\draw[arrow] (output) -- (automated);
\draw[arrow] (output) -- (llmjudge);
\draw[arrow] (output) -- (human);
\draw[arrow] (automated) -- (scores);
\draw[arrow] (llmjudge) -- (scores);
\draw[arrow] (human) -- (scores);
\draw[arrow] (scores) -- (dashboard);

% Feedback loop
\draw[arrow, dashed, bend left=40] (dashboard.south) to node[below, font=\footnotesize, text=gray] {Iterate} (testset.south);

% Labels
\node[above=0.1cm of testset, font=\footnotesize\bfseries, text=green!60!black] {Inputs};
\node[above=0.1cm of llm, font=\footnotesize\bfseries, text=blue!70!black] {System};
\node[above=0.1cm of llmjudge, font=\footnotesize\bfseries, text=orange!70!black] {Evaluation};
\node[above=0.1cm of dashboard, font=\footnotesize\bfseries, text=purple!70!black] {Outputs};

\end{tikzpicture}%
}
\caption{LLM Evaluation Pipeline Overview. Test inputs flow through the LLM application to generate outputs, which are scored by automated metrics, LLM judges, and/or human evaluators. Aggregated scores feed dashboards and inform iterative improvements.}
\label{fig:pipeline-overview}
\end{figure}
